\subsection{The full material derivative at the retained escaping orbit}
\label{sec:material-endpoint}

All coordinates and constants below belong to the original polynomial
\[
\begin{aligned}
F_1&=(1+xy)^3w+y^2(1+xy)(4+3xy),\\
F_2&=y+3x(1+xy)^2w+3xy^2(4+3xy),\\
F_3&=2x-3x^2y-x^3w.
\end{aligned}
\]
Here \(w\) is the identity renaming of the original third spatial coordinate
\(t\). Evolution time is \(\tau\). Put \(J=DF\), \(W_i=J^{-1}e_i\), and
\[
U=2W_1+6F_3W_2,\qquad q_0=(1,-3/2,13/2).
\]
The original determinant is \(\det J=-2\); hence these inverse columns are
defined everywhere. The identities
\[
W_i=-\frac12\nabla F_{i+1}\times\nabla F_{i+2},\qquad
JU=(2,6F_3,0)^{\mathsf T}
\]
use cyclic indices and the original negative determinant.

\subsection{Original time and the positive shorthand}

On the full interval \(-\infty<\tau<1/8\), introduce
\[
z=\sqrt{1-8\tau}>0.
\]
This is a bijection onto \(z>0\), with the exact inverse and derivative
\[
\tau=\frac{1-z^2}{8},\qquad
\frac{dz}{d\tau}=-\frac4z,\qquad
\frac d{d\tau}=-\frac4z\frac d{dz}.
\]
No evolution equation or time coefficient is changed. The initial time
\(\tau=0\) is \(z=1\), and the endpoint \(\tau\uparrow1/8\) is
\(z\downarrow0\).

The trajectory and target material matrix are
\[
\begin{aligned}
\gamma(\tau)&=\left(z^{-1},-\frac32z,\frac{13}{2}z^2\right),\\
B_\tau&=I+6\tau E_{23}
=\begin{pmatrix}
1&0&0\\0&1&\frac34(1-z^2)\\0&0&1
\end{pmatrix}.
\end{aligned}
\]
Here \(E_{23}\) has its only nonzero entry \(1\) in row \(2\), column \(3\).
Direct substitution gives
\[
F(\gamma(\tau))=(-1/4+2\tau,0,0).
\]
Differentiation and \(JU=(2,6F_3,0)^{\mathsf T}\) give
\(\gamma'=U(\gamma)\), with every factor of \(2,6,8\) retained.

\subsection{The complete material derivative and its inverse}

Differentiating the original polynomials and substituting the retained
points gives
\[
\begin{aligned}
J_0:=J(q_0)&=
\begin{pmatrix}
-\frac9{16}&-\frac38&-\frac18\\
\frac38&\frac{25}4&\frac34\\
-\frac{17}2&-3&-1
\end{pmatrix},\\[5pt]
J_\gamma&=
\begin{pmatrix}
-\frac9{16}z^3&-\frac38z&-\frac18\\
\frac38z^2&\frac{25}4&\frac3{4z}\\
-\frac{17}2&-\frac3{z^2}&-\frac1{z^3}
\end{pmatrix}.
\end{aligned}
\]
Both determinants are exactly \(-2\). The full material derivative, with no
coordinate restriction on its input vector, is
\[
A_\tau=J_\gamma^{-1}B_\tau J_0.
\]
Its entire coefficient matrix is
\[
A_\tau=
\begin{pmatrix}
\frac{17z^3-9}{8z^3}&
\frac{3z^3-3}{4z^3}&
\frac{z^3-1}{4z^3}\\[3pt]
\frac{51z^3-24z-27}{16z}&
\frac{9z^3+8z-9}{8z}&
\frac{3z^3-3}{8z}\\[3pt]
-\frac{153}8z^3+\frac92z+\frac{117}8&
-\frac{27}4z^3-3z+\frac{39}4&
\frac{13}4-\frac94z^3
\end{pmatrix}.
\]
Its full inverse is
\[
C_\tau:=A_\tau^{-1}=J_0^{-1}B_\tau^{-1}J_\gamma,
\]
with the entire matrix
\[
C_\tau=
\begin{pmatrix}
\frac{17-9z^3}8&
\frac{3-3z^3}{4z^2}&
\frac{1-z^3}{4z^3}\\[3pt]
\frac{51-24z^2-27z^3}{16}&
\frac{9+8z^2-9z^3}{8z^2}&
\frac{3-3z^3}{8z^3}\\[3pt]
\frac{-153+36z^2+117z^3}8&
\frac{-27-12z^2+39z^3}{4z^2}&
\frac{-9+13z^3}{4z^3}
\end{pmatrix}.
\]
The factored formulas prove both inverse products:
\[
\begin{aligned}
C_\tau A_\tau
&=J_0^{-1}B_\tau^{-1}J_\gamma
 J_\gamma^{-1}B_\tau J_0=I,\\
A_\tau C_\tau&=I.
\end{aligned}
\]
The displayed rational matrices are independently checked entry by entry
against those products. At \(z=1\), they both equal \(I\). Their determinants
satisfy
\[
\begin{aligned}
\det A_\tau
&=(\det J_\gamma)^{-1}\det B_\tau\det J_0\\
&=(-2)^{-1}\cdot1\cdot(-2)=1.
\end{aligned}
\]

\subsection{Full variational equation in the retained time}

Directly differentiating the original polynomial \(U\), before substituting
\(\gamma\), gives
\[
DU(\gamma)=
\begin{pmatrix}
-\frac{27}{2z^2}&-\frac9{z^4}&-\frac3{z^5}\\[3pt]
-\frac{129}4&-\frac{27}{2z^2}&-\frac9{2z^3}\\[3pt]
\frac{423}2z&\frac{93}z&\frac{27}{z^2}
\end{pmatrix}.
\]
Its trace is zero. Applying \(d/d\tau=(-4/z)d/dz\) to the complete
\(A_\tau\) yields
\[
\frac{dA_\tau}{d\tau}
=
\begin{pmatrix}
-\frac{27}{2z^5}&-\frac9{z^5}&-\frac3{z^5}\\[3pt]
-\frac{51}2-\frac{27}{4z^3}&
-9-\frac9{2z^3}&
-3-\frac3{2z^3}\\[3pt]
\frac{459}2z-\frac{18}z&
81z+\frac{12}z&
27z
\end{pmatrix}.
\]
Multiplying the preceding displayed \(DU(\gamma)\) by the displayed
\(A_\tau\) gives this same matrix in all nine entries. Therefore
\[
\frac{dA_\tau}{d\tau}=DU(\gamma(\tau))A_\tau,\qquad
A_0=I,\qquad\det A_\tau=1.
\]
This also confirms that the inverse-Jacobian formula is the actual material
derivative for the original polynomial vector field, not merely a matrix
with the same determinant.

\subsection{The entire transported covector}

Let \(e_3\) mean the third standard initial-coordinate covector, using the
original \((x,y,w)\) basis at \(q_0\). Its transport is
\[
\zeta_\tau=A_\tau^{-\mathsf T}e_3=C_\tau^{\mathsf T}e_3
=
\begin{pmatrix}
\frac{-153+36z^2+117z^3}8\\[3pt]
\frac{-27-12z^2+39z^3}{4z^2}\\[3pt]
\frac{-9+13z^3}{4z^3}
\end{pmatrix}.
\]
Thus its exact Euclidean squared norm is the sum of three squares
\[
\begin{aligned}
|\zeta_\tau|^2
={}&\frac{(-153+36z^2+117z^3)^2}{64}\\
 &+\frac{(-27-12z^2+39z^3)^2}{16z^4}
 +\frac{(-9+13z^3)^2}{16z^6}.
\end{aligned}
\]
In particular
\[
\begin{aligned}
\lim_{z\downarrow0}z^3\zeta_\tau
 &=(0,0,-9/4)^{\mathsf T},\\
\lim_{z\downarrow0}z^6|\zeta_\tau|^2
 &=\frac{81}{16},
\end{aligned}
\]
so
\[
|\zeta_\tau|\sim\frac94z^{-3}
=\frac94(1-8\tau)^{-3/2}.
\]
Here \(f\sim g\) means the ratio tends to \(1\) at the specified endpoint.

There is also an explicit bound on an exact time interval. For
\(0<z\le1/2\), equivalently \(3/32\le\tau<1/8\),
\[
\begin{aligned}
|\zeta_\tau|
&\ge\left|\frac{-9+13z^3}{4z^3}\right|
 =\frac{9-13z^3}{4z^3}\\
&\ge\frac{59}{32}z^{-3}.
\end{aligned}
\]
The last constant follows from \(13z^3\le13/8\); it is a rational
inequality, not a numerical estimate.

\subsection{Every entry of the diffusion tensor}

The full tensor in the initial-coordinate basis is
\[
K_\tau=A_\tau^{-1}A_\tau^{-\mathsf T}
=C_\tau C_\tau^{\mathsf T}.
\]
To state all six independent entries without suppressing any term,
introduce the following seven polynomial abbreviations, local to this
endpoint calculation:
\[
\begin{aligned}
\pi_1(z)&=17-9z^3,&
\pi_2(z)&=1-z^3,\\
\pi_3(z)&=51-24z^2-27z^3,&
\pi_4(z)&=9+8z^2-9z^3,\\
\pi_5(z)&=-153+36z^2+117z^3,&
\pi_6(z)&=-27-12z^2+39z^3,\\
\pi_7(z)&=-9+13z^3.
\end{aligned}
\]
They denote the original numerators. They introduce no identification
with any physical coordinate, density, time-conversion parameter, or
previous scalar field. In the next formulas each \(\pi_j\) means the
displayed \(\pi_j(z)\). The complete symmetric tensor is specified by
\[
\begin{aligned}
K_{11}
 &=\frac{\pi_1^2}{64}
   +\frac{9\pi_2^2}{16z^4}+\frac{\pi_2^2}{16z^6},\\
K_{12}=K_{21}
 &=\frac{\pi_1\pi_3}{128}
   +\frac{3\pi_2\pi_4}{32z^4}+\frac{3\pi_2^2}{32z^6},\\
K_{13}=K_{31}
 &=\frac{\pi_1\pi_5}{64}
   +\frac{3\pi_2\pi_6}{16z^4}+\frac{\pi_2\pi_7}{16z^6},\\
K_{22}
 &=\frac{\pi_3^2}{256}
   +\frac{\pi_4^2}{64z^4}+\frac{9\pi_2^2}{64z^6},\\
K_{23}=K_{32}
 &=\frac{\pi_3\pi_5}{128}
   +\frac{\pi_4\pi_6}{32z^4}+\frac{3\pi_2\pi_7}{32z^6},\\
K_{33}
 &=\frac{\pi_5^2}{64}
   +\frac{\pi_6^2}{16z^4}+\frac{\pi_7^2}{16z^6}.
\end{aligned}
\]
Each equality follows from the corresponding row dot product of the
explicit \(C_\tau\), and the replay also compares all nine entries with
the fully expanded rational tensor. In particular \(K_{33}=|\zeta_\tau|^2\).

For every finite \(\tau<1/8\), \(z>0\), and \(A_\tau,C_\tau\) are
invertible. Therefore, for every nonzero real column \(\xi\),
\[
\xi^{\mathsf T}K_\tau\xi=|C_\tau^{\mathsf T}\xi|^2>0.
\]
This proves strict positive definiteness at every such time. Also
\[
\det K_\tau=(\det C_\tau)^2=1,\qquad K_0=I.
\]
On each compact interval \(0\le\tau\le T'<1/8\), continuity of
\(A_\tau,C_\tau\) gives finite upper bounds on both operator norms.
The exact inequalities
\[
\|A_\tau\|_{\mathrm{op}}^{-2}|\xi|^2
\le \xi^{\mathsf T}K_\tau\xi
\le\|C_\tau\|_{\mathrm{op}}^2|\xi|^2
\]
then give uniform ellipticity on that compact interval.

\subsection{Exact entry powers and all three eigenvalue powers}

Define the following fixed row and column, local to the endpoint
calculation:
\[
\ell=\left(-\frac98,-\frac34,-\frac14\right),\qquad
a_*=\left(\frac14,\frac38,-\frac94\right)^{\mathsf T}.
\]
The symbol \(a_*\) denotes this fixed column and is not a coefficient
coordinate of the fibre cubic. Their squared Euclidean norms are
\[
|\ell|^2=\frac{121}{64},\qquad |a_*|^2=\frac{337}{64}.
\]
The displayed rational expressions give the leading term of every entry
of \(A_\tau\) by rows:
\[
\begin{aligned}
z^3(A_\tau)_{1,*}&\longrightarrow\ell,\\
z(A_\tau)_{2,*}&\longrightarrow\frac32\ell,\\
(A_\tau)_{3,*}&\longrightarrow-13\ell.
\end{aligned}
\]
All entries of these limiting rows are nonzero, so the powers are exact.
The leading terms of every entry of \(C_\tau\), by columns, are
\[
\begin{aligned}
(C_\tau)_{*,1}&\longrightarrow\frac{17}{2}a_*,\\
z^2(C_\tau)_{*,2}&\longrightarrow3a_*,\\
z^3(C_\tau)_{*,3}&\longrightarrow a_*.
\end{aligned}
\]
Again all coefficients are nonzero. Thus the two full matrix limits are
\[
z^3A_\tau\longrightarrow e_1\ell,\qquad
z^3C_\tau\longrightarrow a_*e_3^{\mathsf T}.
\]
Every entry of \(K_\tau\) has exact pole order \(z^{-6}\), with the entire
leading matrix
\[
z^6K_\tau\longrightarrow a_*a_*^{\mathsf T}
=\frac1{64}
\begin{pmatrix}
4&6&-36\\
6&9&-54\\
-36&-54&324
\end{pmatrix}.
\]
The scaled limiting matrix has rank one. This does not change the
previously proved positive definiteness and determinant one of the
unscaled \(K_\tau\) at every finite time.

Let
\[
0<\lambda_1(\tau)\le\lambda_2(\tau)\le\lambda_3(\tau)
\]
be the three eigenvalues of \(K_\tau\). The rank-one limits determine
their full leading behavior without numerical eigenvalue calculations.

Indeed, entrywise convergence of a finite matrix implies convergence in
Frobenius norm, which bounds operator norm. The triangle inequality
\[
\bigl|\|M\|_{\mathrm{op}}-\|N\|_{\mathrm{op}}\bigr|
\le\|M-N\|_{\mathrm{op}}
\le\|M-N\|_{\mathrm{F}}
\]
therefore proves convergence of the operator norms in the two matrix
limits. For any rank-one matrix \(uv^{\mathsf T}\), Cauchy--Schwarz gives
\(\|uv^{\mathsf T}\|_{\mathrm{op}}\le|u||v|\); equality holds on
\(v/|v|\) when \(v\ne0\). Here \(u,v\) in this one elementary norm
identity are arbitrary finite-dimensional columns, not fluid fields.
Applying the identity to the two displayed limits gives
\[
\begin{aligned}
z^3\|A_\tau\|_{\mathrm{op}}
 &\longrightarrow\frac{11}{8},\\
z^3\|C_\tau\|_{\mathrm{op}}
 &\longrightarrow\frac{\sqrt{337}}8.
\end{aligned}
\]
Because \(K_\tau=C_\tau C_\tau^{\mathsf T}\) and
\(K_\tau^{-1}=A_\tau^{\mathsf T}A_\tau\), the largest eigenvalue and the
reciprocal of the smallest are respectively these squared norms.
Consequently
\[
\begin{aligned}
\lambda_1(\tau)
 &\sim\frac{64}{121}z^6
  =\frac{64}{121}(1-8\tau)^3,\\
\lambda_3(\tau)
 &\sim\frac{337}{64}z^{-6}
  =\frac{337}{64}(1-8\tau)^{-3}.
\end{aligned}
\]
Finally \(\det K_\tau=\lambda_1\lambda_2\lambda_3=1\) proves
\[
\lambda_2(\tau)
=\frac1{\lambda_1(\tau)\lambda_3(\tau)}
\longrightarrow\frac{121}{337}.
\]
In particular the middle eigenvalue stays of order one even though every
coordinate entry has a \(z^{-6}\) leading term. The exact condition-number
behavior is
\[
\frac{\lambda_3(\tau)}{\lambda_1(\tau)}
\sim\frac{40777}{4096}z^{-12}
=\frac{40777}{4096}(1-8\tau)^{-6}.
\]

\subsection{Direct exact failure of the two endpoint uniform bounds}

The fixed initial unit covector \(e_3\) has
\[
e_3^{\mathsf T}K_\tau e_3
=|\zeta_\tau|^2
\sim\frac{81}{16}z^{-6}.
\]
Hence there is no finite uniform upper ellipticity bound over
\(0\le\tau<1/8\).

For the lower bound, retain the exact time-dependent unit covector
\[
\xi_\tau=\frac{A_\tau^{\mathsf T}e_1}
 {|A_\tau^{\mathsf T}e_1|}.
\]
It is defined for every finite time because \(A_\tau\) is invertible.
The first row of the full \(A_\tau\) gives
\[
\begin{aligned}
|A_\tau^{\mathsf T}e_1|^2
&=\frac{(17z^3-9)^2+40(z^3-1)^2}{64z^6}\\
&=\frac{329z^6-386z^3+121}{64z^6}.
\end{aligned}
\]
The polynomial \(329z^6-386z^3+121\) is strictly positive: the two
squared factors cannot vanish simultaneously. Since
\(C_\tau^{\mathsf T}A_\tau^{\mathsf T}e_1=e_1\), its exact Rayleigh
quotient is
\[
\begin{aligned}
\xi_\tau^{\mathsf T}K_\tau\xi_\tau
&=\frac{64z^6}{329z^6-386z^3+121}\\
&\sim\frac{64}{121}z^6\longrightarrow0.
\end{aligned}
\]
Thus no positive uniform lower ellipticity constant holds on the full
interval. This is loss of uniform endpoint bounds, with strict positive
definiteness preserved at every finite \(\tau<1/8\).

\subsection{Endpoint replay and scope}

The standalone script \path{scripts/check_material_endpoint.py} retains
the original map, original initial point, original \(\tau\), full
matrices, and all coefficients. Its receipt
\path{checks/material_endpoint_checks.json} contains the expanded entries
as rational functions. All 34 exact checks pass, including the complete
variational equation and all scaled matrix limits. The receipt binds
the present \path{tex/material_endpoint.tex} by its SHA-256 hash.

The result describes the material derivative and tensor along the
specified trajectory. No Lean or numerical inference is used. The
calculation does not assert a classical Navier--Stokes singularity at a
finite original spatial point.

